% Authored. Tables and numbers come from generated/.

\section{The Contribution Base}
\label{sec:contributionbase}

Every Social Security result to this point is internal to the fiscal accounts.
Section~\ref{sec:results} establishes that contributions are the largest positive term
in the recent balance changes, and stops there: an accounting location, not an economic
mechanism. This section supplies the link to something outside the accounts.

\subsection{The base and the ratio}

Contributions are levied on wages, so the natural base is the aggregate wage bill of the
economy, $W_t = N_t \bar{w}_t$, with $N$ employees and $\bar{w}$ the average wage per
employee. Writing the contributions-to-wage-bill ratio as $\tau_t = C_t / W_t$, the
change in contributions decomposes exactly:
\begin{equation}
\Delta C_t = \underbrace{\tau_{t-1} \Delta W_t}_{\text{wage-bill component}}
+ \underbrace{W_{t-1} \Delta \tau_t}_{\text{ratio component}}
+ \underbrace{\Delta W_t \Delta \tau_t}_{\text{interaction}}.
\label{eq:contributionbase}
\end{equation}
The wage-bill component splits again by the same identity, into an employment component
and an average-wage component. Both decompositions are exact. The interaction term is
carried rather than dropped or shared between the factors, because either would make the
decomposition inexact while looking tidier.

We call $\tau$ a ratio and never a rate, deliberately. ``Rate effect'' would invite the
reading that a movement in $\tau$ reflects a change in the statutory levy, and the next
paragraph explains why that reading would be wrong.

The wage bill and the employee count are the Portuguese national accounts compiled by
Statistics Portugal, taken from the harmonised dissemination of those accounts
\citep{ine_wage_bill,ine_employment}. We use wages and salaries (D.11) rather than
compensation of employees (D.1), because D.1 already includes employers' social
contributions and would place part of the numerator inside the denominator.

\paragraph{What $\tau$ is not.} It is not a statutory contribution rate.
National-accounts social contributions received by the Social Security Funds include
imputed contributions and contributions on bases other than employee wages, notably the
self-employed. It is a ratio between two published aggregates, and it moves with
coverage, compliance and composition as well as with legislated rates. It stands at
\EffectiveRateLatest\ in \BaseYear\ against a panel mean of \EffectiveRateMean, both
well below the headline levy, which is itself a sign that the two are not the same
object.

\subsection{What moved contributions}

\begin{table}[htbp]
\centering
\caption{The annual change in Social Security contributions, split into the movement of the wage bill and the movement of the effective ratio, and the wage bill split again into employment and average wages.}
\label{tab:contributionbase}
\small
\fitwidth{%
\begin{tabular}{@{}rrrrrrr@{}}
\toprule
Year & $\Delta$ contributions & From wage bill & From rate & $\Delta$ wage bill & From employment & From average wage \\
\midrule
2018 & 1,166 & 947 & 206 & 3,998 & 1,423 & 2,521 \\
2019 & 1,484 & 1,014 & 443 & 4,225 & 478 & 3,722 \\
2020 & 311 & 28 & 283 & 116 & -1,400 & 1,544 \\
2021 & 1,430 & 1,268 & 152 & 5,071 & 852 & 4,172 \\
2022 & 2,289 & 2,110 & 163 & 8,372 & 2,985 & 5,194 \\
2023 & 2,806 & 2,685 & 107 & 10,572 & 2,000 & 8,383 \\
2024 & 2,557 & 2,056 & 463 & 8,058 & 992 & 6,996 \\
2025 & 2,468 & 1,870 & 560 & 7,196 & 2,253 & 4,842 \\
\bottomrule
\end{tabular}%
}
\par\smallskip
\begin{minipage}{\linewidth}\footnotesize\itshape The interaction terms are carried rather than dropped or shared between the factors, because either would make the decomposition inexact.\end{minipage}
\end{table}


In \BaseYear\ contributions rose by \BaseContributionsChange\ million euro, which
decomposes as
\[
\BaseContributionsChange = \underbrace{\BaseFromWageBill}_{\text{wage bill}}
+ \underbrace{\BaseFromRate}_{\text{ratio}}
+ \underbrace{\BaseRatioInteraction}_{\text{interaction}}.
\]
The wage bill did most of the work, but not all of it. The wage bill itself rose by
\BaseWageBillChange\ million, and that splits the same way:
\[
\BaseWageBillChange = \underbrace{\BaseFromEmployment}_{\text{employees}}
+ \underbrace{\BaseFromAverageWage}_{\text{average wage}}
+ \underbrace{\BaseWageInteraction}_{\text{interaction}}.
\]
The interaction terms are small here, but they are stated rather than absorbed: a paper
that insists the decomposition is exact should not present a two-term version of a
three-term identity.

Figure~\ref{fig:contributionbase} shows both decompositions over the modern period, and
the two crises are where the mechanism is clearest. In the austerity years the wage bill
contracts on both margins at once, employment and wages falling together. In 2020 the
two margins move in opposite directions --- employment falls while average wages rise ---
so the wage bill is almost flat and contributions rise anyway, on the ratio rather than
the base. Recent years are the reverse of austerity: the base grows strongly, and average
wages contribute more than employment.

\begin{figure}[htbp]
\centering
\includegraphics[width=\linewidth]{19_contribution_base.png}
\caption{Contributions and their base, 2001--\PanelEnd, in million euro. Upper panel: the
change in contributions split into a base effect and a rate effect. Lower panel: the
change in the wage bill split into employment and average wages. Both decompositions are
exact; the black line is the identity total in each.}
\label{fig:contributionbase}
\end{figure}

\subsection{A companion regression}

Section~\ref{sec:robustness} sets out why the nominal-GDP specification carries no weight,
and one of its faults was that the regressor was an aggregate that merely correlates with
the balance rather than the base a levy falls on. That fault is repairable, and repairing
it is the point of this specification.

Regressing the annual change in contributions on the annual change in the wage bill over
\WageBillObservations\ adjacent-year pairs gives a slope of \WageBillSlope\ with a HAC
standard error of \WageBillSlopeSe\ and an $R^2$ of \WageBillRSquared. Three things are
worth noting about that.

The slope is close to the average contributions-to-wage-bill ratio of \EffectiveRateMean.
That is consistent with a comparatively stable ratio component over the sample, but it is
not implied by the identity: Equation~\eqref{eq:contributionbase} would deliver
$\beta = \bar{\tau}$ only if $\Delta\tau$ and the interaction term were uncorrelated with
$\Delta W$, which nothing here establishes. The agreement is therefore a descriptive fact
about this sample rather than an arithmetic necessity.

The fit is an order of magnitude tighter than the nominal-GDP regressions, whose $R^2$
values span 0.07 to 0.18. That difference is the substantive argument for the
specification, not a claim about its sophistication.

And no changes are computed across the 1995-to-2000 gap in subsector accounts. That
matters more than it sounds: bridging the gap treats a five-year movement as annual, and
doing so drives the slope to 0.13 and the $R^2$ to 0.40. A single contaminated
observation was enough to destroy the relationship.

\subsection{What this does and does not establish}

The decomposition of Equation~\eqref{eq:contributionbase} is exact and descriptive. It
partitions a recorded movement; it does not show that the wage bill produced it, and it
says nothing about why employment or wages moved. The regression quantifies a
co-movement between two first differences and claims no identification.

What it does establish is that the recent growth in Social Security contributions is
overwhelmingly a base phenomenon rather than a rate phenomenon, and that within the base
it has lately been driven more by wages per employee than by the number of employees.
That is a mechanism in the accounting sense: it names the quantity whose movement the
contributions followed. It is not a causal account of why that quantity moved.
